% =============================================================================
\chapter{Weak Forms and Galerkin Methods}
\label{ch:weakforms}
% =============================================================================

This chapter develops the variational viewpoint behind finite elements. It
shows how strong-form PDE statements become weak statements posed on function
spaces and how Galerkin projections inherit consistency from those weak forms.

\section{Weighted residual statements}
\label{sec:weakforms:weighted_residual}

% Residual orthogonality and the role of test functions

\section{Integration by parts and boundary terms}
\label{sec:weakforms:ibp}

% Reduction of derivative order and emergence of natural boundary conditions

\section{Trial spaces and test spaces}
\label{sec:weakforms:spaces}

% Admissible spaces, lifting functions, and homogeneous reformulations

\section{Galerkin, Petrov-Galerkin, and Ritz formulations}
\label{sec:weakforms:galerkin}

% Bubnov-Galerkin, Petrov-Galerkin, and energy minimization viewpoints

\section{Consistency, stability, and orthogonality}
\label{sec:weakforms:consistency}

% Galerkin orthogonality, residual consistency, and stability requirements

\section{Abstract variational problems}
\label{sec:weakforms:abstract}

% Bilinear forms, linear functionals, and operator notation
